\section{Discussion}
\label{sec:discussion}

\subsection{What the account requires}

Read against the models, William's account describes a narrow set of conditions.
A flight of a furlong from a tower of this height needs a strong southwest wind, a wing
near the better end of what cloth and timber allow, and a track that stays over the
slope.
Flying straight out over the valley, the route \textcite{woosnam1986eilmer} proposed, does
not reach \SI{200}{\meter} in any sampled case, because the headwind slows the craft over
the ground and the lift ends at the foot of the slope.
Flying along the slope reaches it in \SI{\nReachSlope}{\percent} of cases.
Southwest winds of \SI{8}{\meter\per\second} or more blow for about \nSWEightHours\ hours a
year, so a man who waited for one would not have waited long.
William's phrase that Eilmer ``collected the breeze'' fits a launch timed to such a wind.

Without a tail the craft diverges in pitch in a fraction of a second, but the doubling time
is close to human reaction time rather than far below it.
A pilot who moves his body could therefore hold the attitude for a while by continual
correction, and a flight of tens of seconds followed by a loss of control is consistent
with the model.
The account names the trigger: the violence of the wind and the swirling of the air.
A gust large enough to need a correction faster than his reactions would end the flight in
under a second (\cref{fig:stability}a).
The models do not show that he lost control, only that he would have been unlikely to keep
it and that a tail would have removed the need to.

The broken legs are consistent with a loss of control but do not prove one.
A controlled glide into the ground without a flare arrives too slowly to break the legs in
this model, so the injury requires a fast, steep arrival.
A loss of pitch control gives that, but so would a stall at height or a dive into the
ground that is not caused by the missing tail.

\subsection{Which results depend on the assumptions}

Three results follow from measurement or from simple bounds and would survive large
changes to the model.
The slope faces the prevailing wind, and strong southwest winds are measured to occur
(\cref{fig:site,fig:wind}).
The launch dive costs \((V^2-v_0^2)/2g\) of height whatever the wing, so a still-air flight
of a furlong is not possible from this tower.
A tailless wing is unstable in pitch whenever the CG is aft of about \num{\nCGonset}\,\(c\),
and for a man fixed to the wing at his hands and feet the CG is near mid-chord.
This last result holds if the sampled CG range brackets the real posture;
\cref{fig:stability}b shows it would fail only if the pilot's mass centre lay ahead of
\num{\nCGonset}\,\(c\), within the front of the chord near his hands.

The fraction of cases reaching \SI{200}{\meter} depends on the ranges in Supplementary
\cref{tab:priors}, and should be read as the share of the sampled ranges that works, not as
a probability that the flight happened.
It is most sensitive to the glide ratio and the strength and length of the ridge lift.
The glide ratio rests on \cref{eq:polar} with bounds on aspect ratio, span efficiency, and
drag, and every wing is flown at its best glide, which favours the account.
The lift rests on a first-order model over a smoothed \SI{6}{\meter} terrain grid; the real
slope is steeper than the grid shows, which would raise the lift, while turbulence and
separation behind the spur edge would lower it.
The point-mass model holds \(C_L\) fixed, so it does not let the pilot trade height for speed
more efficiently than a stable glider would, and the pitch analysis treats the longitudinal
motion alone.
A large-eddy simulation over LIDAR terrain and a six-degree-of-freedom model with a pilot
model would narrow these margins.

\subsection{Limits of the comparison}

The comparison with the text assumes that a furlong means about \SI{200}{\meter} of ground,
that the tower stood at or near the abbey, and that William reports what Eilmer said.
A shorter flight would be easier to explain; a tower on lower ground would be harder.
None of the models tests whether the structure could carry the flight loads, which would be
a separate failure mode.
